% ============================================================
%  APPENDIX A --- Installation
% ============================================================
\chapter{Installation}
\label{app:installation}

This appendix covers every supported deployment scenario for FairWindSK,
from a first-time source build on a developer laptop to a fully automatic
kiosk boot on a dedicated helm display.
Each scenario lists hardware or platform requirements, exact commands,
and post-install verification steps.

\begin{notebox}
All scenarios assume a running Signal~K server somewhere on the boat's
Wi-Fi network.
FairWindSK is a display shell, not a data server.
Install and configure Signal~K before installing FairWindSK.
\end{notebox}

% ── A.1  Prerequisites ───────────────────────────────────────
\section{Prerequisites}
\label{app:prerequisites}

\subsection{Common prerequisites for all desktop platforms}

\begin{itemize}
  \item CMake 3.16 or newer
  \item A C++17-capable compiler supported by the chosen Qt~6 kit
  \item Git
  \item Internet access during the first build
        (CMake downloads \code{nlohmann/json}, \code{QtZeroConf}, and \code{QHotkey}
        automatically; subsequent builds use the cached copies)
\end{itemize}

\subsection{Required Qt 6 modules --- desktop}

\rowcolors{2}{LtGray}{white}
\begin{longtable}{L{4.5cm}L{9.2cm}}
\toprule
\rowcolor{Navy}
\color{white}\sffamily\bfseries Module &
\color{white}\sffamily\bfseries Purpose \\
\midrule
\code{Core}            & Qt base classes and event loop. \\
\code{Gui}             & 2D graphics and font handling. \\
\code{Widgets}         & Widget toolkit used throughout the shell. \\
\code{Concurrent}      & Thread pool for background tasks. \\
\code{Network}         & HTTP/REST client and TLS. \\
\code{WebSockets}      & Signal~K WebSocket stream. \\
\code{Xml}             & Configuration and resource parsing. \\
\code{Svg}             & OpenBridge icon rendering. \\
\code{Positioning}     & GPS location APIs (desktop fallback). \\
\code{WebEngineWidgets}& Embedded browser for hosted Signal~K web apps. \\
\code{VirtualKeyboard} & On-screen keyboard for touchscreen deployments. \\
\code{PrintSupport}    & PDF generation (crash reports). \\
\bottomrule
\end{longtable}

\subsection{Required Qt 6 modules --- Android and iOS}

Mobile builds replace \code{WebEngineWidgets} with a lighter web layer and
drop desktop-only dependencies:

\rowcolors{2}{LtGray}{white}
\begin{longtable}{L{4.5cm}L{9.2cm}}
\toprule
\rowcolor{Navy}
\color{white}\sffamily\bfseries Module &
\color{white}\sffamily\bfseries Replaces or adds \\
\midrule
\code{Qml}         & QML runtime for WebView hosting. \\
\code{Quick}       & Qt Quick scene graph. \\
\code{QuickWidgets}& \code{QQuickWidget} bridge between widget and QML worlds. \\
\code{WebView}     & Replaces \code{WebEngineWidgets}. Uses the platform's native browser engine. \\
\bottomrule
\end{longtable}

Mobile builds do \emph{not} include \code{QtZeroConf}, \code{QHotkey},
\code{PrintSupport}, or \code{WebEngineWidgets}.

% ── A.2  Generic desktop build ───────────────────────────────
\section{Generic Desktop Build (all platforms)}
\label{app:generic}

This is the baseline procedure that applies to every desktop platform
before platform-specific steps are added.

\begin{enumerate}
  \item Clone the repository:
\begin{lstlisting}
git clone https://github.com/OpenFairWind/FairWindSK.git
cd FairWindSK
\end{lstlisting}

  \item Configure:
\begin{lstlisting}
cmake -S . -B build
\end{lstlisting}

  \item Build (parallel across all CPU cores):
\begin{lstlisting}
cmake --build build --parallel
\end{lstlisting}

  \item Install to the default prefix:
\begin{lstlisting}
cmake --install build
\end{lstlisting}

  \item Override the install prefix if needed:
\begin{lstlisting}
cmake -S . -B build -DCMAKE_INSTALL_PREFIX=/desired/prefix
cmake --build build --parallel
cmake --install build
\end{lstlisting}

  \item If Qt is installed outside the default search path, pass
        \code{-DCMAKE\_PREFIX\_PATH} during configure:
\begin{lstlisting}
cmake -S . -B build -DCMAKE_PREFIX_PATH=/path/to/Qt/6.x.x/gcc_64
\end{lstlisting}
\end{enumerate}

\begin{tipbox}
If internet access is restricted during the build, install
\code{nlohmann\_json} from your package manager first.
CMake will use the system header instead of downloading it.
On Debian/Ubuntu: \code{sudo apt install nlohmann-json3-dev}.
On macOS with Homebrew: \code{brew install nlohmann-json}.
\end{tipbox}

% ── A.3  macOS ───────────────────────────────────────────────
\section{macOS Installation}
\label{app:macos}

\subsection{Requirements}

\begin{itemize}
  \item macOS 12 Monterey or newer (recommended)
  \item Xcode command line tools: \code{xcode-select --install}
  \item Qt~6 desktop kit (via Qt Online Installer or Homebrew)
  \item CMake and Ninja (via Homebrew or Qt Maintenance Tool)
\end{itemize}

\subsection{Install with Homebrew}

\begin{lstlisting}
brew install qt cmake ninja nlohmann-json
cmake -S . -B build -G Ninja \
  -DCMAKE_PREFIX_PATH="$(brew --prefix qt)"
cmake --build build --parallel
\end{lstlisting}

\subsection{Run from the build tree}

\begin{lstlisting}
open build/FairWindSK.app
\end{lstlisting}

\subsection{Install to the system}

\begin{lstlisting}
cmake --install build
# Installs to /usr/local by default:
open /usr/local/FairWindSK.app
\end{lstlisting}

\subsection{Redistribute (deploy Qt runtime)}

For distribution outside the build tree, use Apple's deployment tool:

\begin{lstlisting}
macdeployqt build/FairWindSK.app
\end{lstlisting}

This bundles all Qt frameworks inside the \code{.app} so the app runs
on machines without Qt installed.

\subsection{Post-install verification}

\begin{enumerate}
  \item Launch FairWindSK.
  \item Open \ui{Settings~> Connection}.
        Enter your Signal~K server address and tap \key{Connect}.
  \item Confirm the status changes to \textit{Live} and apps appear on the launcher.
\end{enumerate}

% ── A.4  Linux desktop ───────────────────────────────────────
\section{Linux Desktop Installation}
\label{app:linux}

\subsection{Requirements (Debian / Ubuntu)}

\begin{lstlisting}
sudo apt update
sudo apt install \
  build-essential cmake ninja-build git \
  nlohmann-json3-dev \
  qt6-base-dev qt6-base-dev-tools \
  qt6-webengine-dev qt6-webengine-dev-tools \
  qt6-websockets-dev qt6-positioning-dev \
  libqt6svg6-dev qt6-virtualkeyboard-dev \
  qt6-base-private-dev libqt6printsupport6 \
  libavahi-compat-libdnssd-dev libnss-mdns avahi-utils
\end{lstlisting}

\subsection{Build and run}

\begin{lstlisting}
cmake -S . -B build -G Ninja
cmake --build build --parallel
./build/FairWindSK
\end{lstlisting}

\subsection{System install}

\begin{lstlisting}
sudo cmake --install build
\end{lstlisting}

\code{cmake --install} on a Linux desktop target:

\begin{itemize}
  \item Installs the \code{FairWindSK} binary to
        \code{\$\{CMAKE\_INSTALL\_BINDIR\}} (default \code{/usr/local/bin}).
  \item Installs the bundled icon directory alongside the binary.
  \item Installs the desktop-only helper libraries to
        \code{\$\{CMAKE\_INSTALL\_LIBDIR\}}.
  \item Installs the \code{fairwindsk.desktop} launcher file to the standard
        XDG applications directory.
  \item Installs the \code{fairwindsk.png} menu icon to
        \code{/usr/local/share/icons/hicolor/256x256/apps/}.
\end{itemize}

After installation, FairWindSK appears in the applications menu.
Run it from the menu or with:

\begin{lstlisting}
FairWindSK
\end{lstlisting}

\subsection{Packaging}

Include the Qt runtime and the external dependencies built under
\code{build/external/lib} in your package.
The build-tree rpath means the binary runs from \code{build/} without
\code{LD\_LIBRARY\_PATH}, but a packaged install needs those libraries in
the system library path or bundled beside the binary.

% ── A.5  Raspberry Pi OS --- nav station ───────────────────────
\section{Raspberry Pi OS --- Navigation Station}
\label{app:rpi-station}

This scenario installs FairWindSK on a Raspberry~Pi used as a normal
desktop navigation station --- screen, keyboard, and mouse attached,
running the Raspberry~Pi~OS desktop.

\subsection{Requirements}

Raspberry Pi~4 or~5 with Raspberry~Pi~OS Bookworm (64-bit recommended),
and a stable internet connection for the initial package download.

\subsection{Install build dependencies}

\begin{lstlisting}
sudo apt update
sudo apt install \
  build-essential cmake ninja-build git \
  nlohmann-json3-dev \
  qml6-module-qt-labs-folderlistmodel \
  qml6-module-qtquick-window \
  qml6-module-qtquick-layouts \
  qml6-module-qtqml-workerscript \
  libnss-mdns avahi-utils \
  libavahi-compat-libdnssd-dev libxkbcommon-dev \
  qt6-base-dev qt6-base-dev-tools \
  qt6-webengine-dev qt6-webengine-dev-tools \
  qt6-websockets-dev qt6-positioning-dev \
  libqt6svg6-dev \
  qt6-virtualkeyboard-dev \
  libqt6virtualkeyboard6 \
  qt6-virtualkeyboard-plugin
\end{lstlisting}

\begin{notebox}
If the build fails with a missing
\code{Qt6VirtualKeyboardConfigVersionImpl.cmake}, apply the known
workaround:
\begin{lstlisting}
sudo touch /usr/lib/aarch64-linux-gnu/cmake/Qt6VirtualKeyboard/\
Qt6VirtualKeyboardConfigVersionImpl.cmake
\end{lstlisting}
Then re-run \code{cmake}.
\end{notebox}

\subsection{Build and install}

\begin{lstlisting}
git clone https://github.com/OpenFairWind/FairWindSK.git
cd FairWindSK
cmake -S . -B build -G Ninja
cmake --build build --parallel
sudo cmake --install build
\end{lstlisting}

FairWindSK appears in the application menu under
\textit{Navigation} or \textit{Internet}.

\subsection{OpenPlotter integration}

If OpenPlotter is detected during \code{cmake --install}, FairWindSK
also copies its launcher icon into the OpenPlotter menu directories.
No extra steps are required.

% ── A.6  Raspberry Pi OS --- kiosk helm display ────────────────
\section{Raspberry Pi OS --- Kiosk Helm Display}
\label{app:rpi-kiosk}

This is the recommended scenario for a dedicated helm display:
a Raspberry~Pi~4 or~5 with a touchscreen, running FairWindSK automatically
at boot with no keyboard and no desktop menu visible.

\subsection{Overview}

\begin{enumerate}
  \item Install FairWindSK as in Section~\ref{app:rpi-station}.
  \item Configure FairWindSK for full-screen kiosk mode.
  \item Install the \code{fairwindsk-startup} helper to manage autostart.
  \item (Optional) Configure the desktop panel to auto-hide.
  \item Reboot and verify.
\end{enumerate}

\subsection{Step 1 --- Build and install}

Follow Section~\ref{app:rpi-station}.
On Raspberry~Pi~OS, \code{cmake --install} automatically installs the
\code{fairwindsk-startup.desktop} autostart entry to
\code{/etc/xdg/autostart/}.
FairWindSK will start automatically on the next login.

\subsection{Step 2 --- Configure full-screen kiosk mode}

Open \ui{Settings~> Main} and set \key{Window Mode} to
\textbf{Full Screen}.
This requests a frameless, always-on-top window that covers the entire
screen with no title bar or window decorations.

Alternatively, edit \code{fairwindsk.json} directly:

\begin{lstlisting}
{
  "main": {
    "windowMode": "fullscreen"
  }
}
\end{lstlisting}

\subsection{Step 3 --- Install the startup helper}

The \code{fairwindsk-startup} Python script manages the FairWindSK
process lifecycle.
It is installed by \code{cmake --install} alongside the binary.

What \code{fairwindsk-startup} does at each boot:

\begin{enumerate}
  \item Reads \code{fairwindsk.ini} to find the configuration path.
  \item Reads \code{fairwindsk.json} to check the \code{virtualKeyboard}
        flag and sets \code{QT\_IM\_MODULE=qtvirtualkeyboard} accordingly.
  \item Waits up to 45 seconds for the Signal~K web-app catalogue
        endpoint to respond (\code{/signalk/v1/apps/list} or the legacy
        \code{/skServer/webapps}).
        This ensures plugin icons and artwork are available before
        FairWindSK starts.
  \item Launches FairWindSK.
  \item If FairWindSK exits with code~1 (the \textit{Restart} button in
        \ui{Settings~> System}), the helper relaunches it immediately.
  \item If FairWindSK exits with code~0 (the \textit{Quit} button or
        a clean shutdown), the helper exits cleanly.
\end{enumerate}

\begin{warningbox}
Use \key{Restart FairWindSK} in \ui{Settings~> System} (exit code~1)
to relaunch after configuration changes.
\key{Quit FairWindSK} (exit code~0) leaves the display blank
without a restart.
\end{warningbox}

\subsection{Step 4 --- Hide the desktop panel (labwc / Bookworm)}

On Raspberry~Pi~OS Bookworm with the labwc compositor, the taskbar
panel can overlay the full-screen FairWindSK window.
To prevent this, configure the panel to auto-hide:

\begin{lstlisting}
# Edit the panel configuration file
nano ~/.config/wf-panel-pi.ini
\end{lstlisting}

Add or change:

\begin{lstlisting}
[panel]
autohide = true
\end{lstlisting}

Save the file and reboot.

\subsection{Step 5 --- Enable virtual keyboard (touchscreen only)}

If the helm display has no physical keyboard, enable the Qt virtual
keyboard:

\begin{enumerate}
  \item Open \ui{Settings~> Main}.
  \item Tick \key{Virtual Keyboard}.
  \item Tap \key{Restart FairWindSK} in \ui{Settings~> System}.
\end{enumerate}

The startup helper reads the setting at each boot and sets the
environment variable automatically.

\subsection{Step 6 --- Verify the kiosk boot}

\begin{enumerate}
  \item Reboot the Raspberry~Pi.
  \item Verify FairWindSK starts automatically (allow up to 60 seconds
        for Signal~K and FairWindSK to start together).
  \item Confirm the display is full-screen with no title bar.
  \item Confirm the status shows \textit{Live}.
  \item Confirm apps appear on the launcher.
\end{enumerate}

\begin{tipbox}
On a Raspberry~Pi~4 with a class~A2 microSD card, FairWindSK
typically becomes interactive in under 30 seconds from power-on,
assuming Signal~K starts quickly.
Use a fast card (\textgreater{}90~MB/s read) to minimise boot time.
\end{tipbox}

% ── A.7  Windows ─────────────────────────────────────────────
\section{Windows Installation}
\label{app:windows}

\subsection{Requirements}

\begin{itemize}
  \item Windows~10 or~11 (64-bit)
  \item Qt~6 desktop kit for MSVC
        (installed via the Qt Online Installer)
  \item Visual Studio Build Tools 2022 or Visual Studio Community 2022
  \item CMake and Ninja
        (bundled with Qt or installed separately)
\end{itemize}

\subsection{Build and install}

Open a Developer Command Prompt with the MSVC environment loaded,
then:

\begin{lstlisting}
git clone https://github.com/OpenFairWind/FairWindSK.git
cd FairWindSK
cmake -S . -B build -G Ninja ^
  -DCMAKE_PREFIX_PATH="C:\Qt\6.x.x\msvc2022_64"
cmake --build build --parallel
cmake --install build
\end{lstlisting}

\subsection{Deploy the Qt runtime}

After a successful build, deploy the Qt runtime DLLs beside the
executable:

\begin{lstlisting}
windeployqt build\FairWindSK.exe
\end{lstlisting}

\code{cmake --install} installs the executable, the bundled icon
directory under \code{bin\textbackslash icons}, and the desktop-only
helper DLLs into the same \code{bin} directory.

\subsection{Run}

\begin{lstlisting}
build\FairWindSK.exe
\end{lstlisting}

\begin{notebox}
\code{file://} launcher entries are retained in configuration for
compatibility but are blocked by the single-window model on Windows.
Use \code{http://} or \code{https://} URLs for all applications.
\end{notebox}

% ── A.8  Android ─────────────────────────────────────────────
\section{Android Installation}
\label{app:android}

Android builds use the Qt WebView mobile implementation.
Qt Creator is the recommended workflow for Android because it manages
the SDK, NDK, APK signing, and device deployment.

\subsection{Requirements}

\begin{itemize}
  \item A macOS, Linux, or Windows build host
  \item Qt~6 with the Android kit for your target ABI
        (arm64-v8a is recommended for modern devices)
  \item Android Studio (for the SDK, NDK, and emulator)
  \item Java JDK (bundled with Android Studio or installed separately)
  \item CMake and Ninja
\end{itemize}

\subsection{Step-by-step build guide}

\begin{enumerate}
  \item Open Android Studio and install:
        an Android SDK platform, SDK Platform-Tools, SDK Build-Tools,
        and the NDK version recommended by your Qt release.
  \item In the Qt Maintenance Tool, install the Android Qt components
        that match your host and target ABI.
  \item In Qt Creator, go to
        \ui{Preferences~> Devices~> Android} and verify that the SDK,
        NDK, JDK, and platform paths are all detected correctly.
  \item Clone the repository:
\begin{lstlisting}
git clone https://github.com/OpenFairWind/FairWindSK.git
\end{lstlisting}
  \item Open \code{FairWindSK/CMakeLists.txt} in Qt Creator.
  \item Select an Android kit (e.g.\ \textit{Android Qt~6.x.x Clang arm64-v8a}).
  \item Let Qt Creator configure the project.
        The mobile configuration uses \code{Qt::WebView},
        \code{Qt::Quick}, and \code{Qt::QuickWidgets},
        and skips desktop-only dependencies.
  \item Build with \ui{Build~> Build Project}.
  \item Connect an Android device with developer mode enabled,
        or start an Android emulator.
  \item Deploy with \ui{Build~> Run}.
\end{enumerate}

\subsection{Command-line configure (advanced)}

\begin{lstlisting}
cmake -S . -B build-android -G Ninja \
  -DCMAKE_TOOLCHAIN_FILE="$ANDROID_NDK/build/cmake/android.toolchain.cmake" \
  -DANDROID_ABI=arm64-v8a \
  -DANDROID_PLATFORM=android-24 \
  -DCMAKE_PREFIX_PATH="/path/to/Qt/6.x.x/android_arm64_v8a"
cmake --build build-android --parallel
\end{lstlisting}

Qt Creator is still recommended for packaging and device deployment.

\subsection{Android feature differences}

\rowcolors{2}{LtGray}{white}
\begin{longtable}{L{5.2cm}L{8.5cm}}
\toprule
\rowcolor{Navy}
\color{white}\sffamily\bfseries Feature &
\color{white}\sffamily\bfseries Android status \\
\midrule
Hosted web apps          & Supported via Qt WebView. \\
Signal~K stream          & Supported (WebSockets). \\
mDNS server discovery    & Not available (disabled on Android). \\
Global hotkey SHIFT+TAB  & Not available (desktop-only, uses QHotkey). \\
Virtual keyboard         & Uses Android system keyboard; Qt Virtual Keyboard setting has no effect. \\
File printing / PDF      & Not available (PrintSupport not included). \\
QWebEngine cookies       & Not used; mobile path uses a separate mechanism. \\
\code{file://} apps      & Blocked by the single-window model on all targets. \\
\bottomrule
\end{longtable}

% ── A.9  iOS and iPadOS ──────────────────────────────────────
\section{iOS and iPadOS Installation}
\label{app:ios}

iOS and iPadOS builds use the same mobile web implementation as Android.
Qt Creator with an Apple developer account is required for device deployment.

\subsection{Requirements}

\begin{itemize}
  \item A macOS build host
  \item Xcode and Xcode command line tools
  \item Qt~6 with the iOS kit
  \item CMake and Ninja
  \item An Apple developer account and provisioning profile
        (for device deployment; not needed for the Simulator)
\end{itemize}

\subsection{Step-by-step build guide}

\begin{enumerate}
  \item Open Xcode at least once and accept the license agreement.
  \item In the Qt Maintenance Tool, install the iOS Qt components
        that match your desktop Qt release.
  \item Clone the repository:
\begin{lstlisting}
git clone https://github.com/OpenFairWind/FairWindSK.git
\end{lstlisting}
  \item Open \code{FairWindSK/CMakeLists.txt} in Qt Creator.
  \item Select an iOS kit (e.g.\ \textit{Qt~6.x.x for iOS}).
  \item Build with \ui{Build~> Build Project}.
  \item To run on the Simulator: choose an iOS Simulator run target
        and select \ui{Build~> Run}.
  \item To run on a device: connect the iPhone or iPad, select the
        physical device in the run target chooser, verify signing settings
        when prompted, and run from Qt Creator.
\end{enumerate}

\subsection{Command-line configure (advanced)}

\begin{lstlisting}
cmake -S . -B build-ios -G Ninja \
  -DCMAKE_SYSTEM_NAME=iOS \
  -DCMAKE_OSX_SYSROOT=iphoneos \
  -DCMAKE_OSX_ARCHITECTURES=arm64 \
  -DCMAKE_PREFIX_PATH="/path/to/Qt/6.x.x/ios"
cmake --build build-ios --parallel
\end{lstlisting}

Use \code{iphonesimulator} and \code{x86\_64} (or \code{arm64} on
Apple Silicon hosts) for Simulator builds.

\subsection{iOS feature differences}

\rowcolors{2}{LtGray}{white}
\begin{longtable}{L{5.2cm}L{8.5cm}}
\toprule
\rowcolor{Navy}
\color{white}\sffamily\bfseries Feature &
\color{white}\sffamily\bfseries iOS / iPadOS status \\
\midrule
Hosted web apps          & Supported via Qt WebView (WKWebView). \\
Signal~K stream          & Supported (WebSockets). \\
mDNS server discovery    & Not available. Enter the server URL manually. \\
Global hotkey SHIFT+TAB  & Not available. \\
Virtual keyboard         & Uses iOS system keyboard. \\
File printing / PDF      & Not available. \\
\code{file://} apps      & Blocked by the single-window model on all targets. \\
\bottomrule
\end{longtable}

% ── A.10  Signal K Server on Docker ──────────────────────────
\section{Signal K Server on Docker}
\label{app:signalk-docker}

If you do not already have a Signal~K server, Docker is the fastest
way to get one running.
This section documents the Docker deployment from the
FairWindSK project repository.

\subsection{Docker image}

\begin{lstlisting}
cr.signalk.io/signalk/signalk-server:latest
\end{lstlisting}

Supported architectures: \code{linux/amd64}, \code{linux/arm/v7},
\code{linux/arm64} (covers all Raspberry~Pi~4 and~5 configurations).

\subsection{Quick start with demo NMEA data}

\begin{lstlisting}
mkdir $HOME/.signalk
docker run --init -it --rm \
  --name signalk-server \
  --publish 3000:3000 \
  -v $HOME/.signalk:/home/node/.signalk \
  cr.signalk.io/signalk/signalk-server \
  --sample-nmea0183-data
\end{lstlisting}

This runs Signal~K on port~3000 with simulated NMEA~0183 data.
Use it to test FairWindSK without real instruments.

\subsection{Production deployment}

\begin{lstlisting}
docker run -d --init \
  --name signalk-server \
  -p 3000:3000 \
  -v $HOME/.signalk:/home/node/.signalk \
  cr.signalk.io/signalk/signalk-server
\end{lstlisting}

\begin{itemize}
  \item The \code{-d} flag runs the container as a background daemon.
  \item The \code{-v} flag mounts the configuration directory so
        Signal~K settings, plugins, and data persist across container
        restarts.
  \item On first access, the Signal~K admin web UI at
        \code{http://localhost:3000} prompts you to create admin credentials.
\end{itemize}

\subsection{Updating Signal K}

\begin{lstlisting}
docker pull cr.signalk.io/signalk/signalk-server:latest
docker stop signalk-server
docker rm signalk-server
# Re-run the production docker run command above
\end{lstlisting}

\subsection{Directory structure inside the container}

\rowcolors{2}{LtGray}{white}
\begin{longtable}{L{4.8cm}L{8.9cm}}
\toprule
\rowcolor{Navy}
\color{white}\sffamily\bfseries Path &
\color{white}\sffamily\bfseries Contents \\
\midrule
\code{/home/node/signalk}  & Signal~K server files (read-only inside container). \\
\code{/home/node/.signalk} & Settings files, plugins, and data.
                              \textbf{Mount this to the host.} \\
\bottomrule
\end{longtable}

% ── A.11  First-Run Configuration ────────────────────────────
\section{First-Run Configuration}
\label{app:firstrun}

After installation on any platform, follow these steps on the first launch.

\subsection{Configuration file locations}

FairWindSK stores its configuration in two files in the per-user
Qt configuration directory:

\rowcolors{2}{LtGray}{white}
\begin{longtable}{L{3.6cm}L{10.1cm}}
\toprule
\rowcolor{Navy}
\color{white}\sffamily\bfseries File &
\color{white}\sffamily\bfseries Contents \\
\midrule
\code{fairwindsk.json}  &
  All user-visible settings: connection URL, bar layouts, app catalogue,
  widget definitions, comfort presets, units, and Signal~K path mappings.
  Human-readable JSON.
  Safe to back up, transfer, and version-control. \\
\code{fairwindsk.ini}   &
  Qt-managed settings: path to \code{fairwindsk.json}, the debug flag,
  and the authentication token.
  The token is kept here intentionally so it is never included in a
  configuration export. \\
\bottomrule
\end{longtable}

On first run, if \code{fairwindsk.json} does not exist, FairWindSK
seeds it from the bundled defaults in
\code{resources/json/configuration.json}.

\subsection{Connect to Signal K}

\begin{enumerate}
  \item Open \ui{Settings~> Connection}.
  \item Enter the Signal~K server address,
        e.g.\ \code{http://192.168.1.100:3000}.
        If you are running Signal~K on the same device,
        use \code{http://localhost:3000}.
  \item Tap \key{Connect}.
  \item Wait for the status to change to \textit{Live}.
  \item If the server requires authentication,
        tap \key{Request Token} and approve the request in the
        Signal~K admin interface.
\end{enumerate}

\subsection{Authenticate (optional)}

If the Signal~K server has security enabled:

\begin{enumerate}
  \item Tap \key{Request Token} in \ui{Settings~> Connection}.
  \item FairWindSK opens the Signal~K admin interface in the
        Bottom~Bar drawer, navigated to
        \ui{Security~> Access Requests}.
  \item Approve the request as an admin user.
  \item FairWindSK stores the token automatically and reconnects.
\end{enumerate}

\subsection{Essential first-time settings}

\rowcolors{2}{LtGray}{white}
\begin{longtable}{L{4.0cm}L{9.7cm}}
\toprule
\rowcolor{Navy}
\color{white}\sffamily\bfseries Setting &
\color{white}\sffamily\bfseries Where and what to set \\
\midrule
Units &
  \ui{Settings~> Units}.
  Choose knots (speed), nautical miles (distance), and metres or feet (depth). \\
Coordinate format &
  \ui{Settings~> Main}.
  Choose DD°MM.MMM' for standard nautical use. \\
Window mode &
  \ui{Settings~> Main}.
  Choose Full~Screen for a dedicated helm display;
  Windowed or Centred for a navigation station. \\
Comfort preset &
  \ui{Settings~> Comfort}.
  Set Day for initial testing; change to Night for the night watch. \\
Launcher apps &
  \ui{Settings~> Applications}.
  Arrange your chart plotter, instruments, and other apps on page~1. \\
Virtual keyboard &
  \ui{Settings~> Main} (touchscreen only).
  Tick, then restart FairWindSK. \\
\bottomrule
\end{longtable}

% ── A.12  Configuration file reference ───────────────────────
\section{Key Configuration Fields}
\label{app:config}

Advanced users and system integrators can edit \code{fairwindsk.json}
directly.
This section covers the most useful fields.

\subsection{Connection}

\begin{lstlisting}
{
  "connection": {
    "server": "http://your-signalk-host:3000"
  }
}
\end{lstlisting}

\subsection{Window mode}

\begin{lstlisting}
{
  "main": {
    "windowMode": "fullscreen",
    "language": "en",
    "virtualKeyboard": false
  }
}
\end{lstlisting}

Valid values for \code{windowMode}: \code{windowed}, \code{centered},
\code{maximized}, \code{fullscreen}.
Valid values for \code{language}: \code{system}, \code{en}, \code{it}.

\subsection{Units}

\begin{lstlisting}
{
  "units": {
    "vesselSpeed": "kn",
    "windSpeed":   "kn",
    "distance":    "nm",
    "depth":       "mt"
  }
}
\end{lstlisting}

\subsection{Signal K path overrides}

\begin{lstlisting}
{
  "signalk": {
    "sog": "navigation.speedOverGround",
    "dpt": "environment.depth.belowTransducer",
    "notifications.pob": "notifications.mob"
  }
}
\end{lstlisting}

\subsection{Diagnostics}

\begin{lstlisting}
{
  "diagnostics": {
    "logLevel":          "off",
    "persistentLogs":    true,
    "interactionHistory": true,
    "email": "your-email@example.com"
  }
}
\end{lstlisting}

Valid log levels: \code{off}, \code{critical}, \code{warning},
\code{info}, \code{debug}, \code{full}.

% ── A.13  Troubleshooting installation ───────────────────────
\section{Troubleshooting Installation Problems}
\label{app:troubleshooting-install}

\subsection{FairWindSK does not start after installation}

\begin{enumerate}
  \item Run the binary from a terminal to see error output:
        \code{FairWindSK} (Linux/macOS) or
        \code{build\textbackslash FairWindSK.exe} (Windows).
  \item Check that Qt WebEngine libraries are available.
        On Linux, run \code{ldd build/FairWindSK | grep "not found"}.
  \item On Raspberry~Pi, check for the virtual keyboard CMake workaround
        (Section~\ref{app:rpi-station}).
\end{enumerate}

\subsection{No applications appear on the launcher}

\begin{enumerate}
  \item Verify the Signal~K server address in
        \ui{Settings~> Connection}.
  \item Confirm the server exposes
        \code{/signalk/v1/apps/list} or \code{/skServer/webapps}.
  \item Check the Signal~K App Store --- at least one web app must be
        installed on the server for it to appear in FairWindSK.
  \item Enable debug logging: add \code{debug=true} to
        \code{fairwindsk.ini} and relaunch.
        Logs include connection attempts and app discovery details.
\end{enumerate}

\subsection{Build fails with missing Qt module}

\begin{enumerate}
  \item Verify the Qt version: \code{qmake6 --version} or
        \code{cmake --find-package -DNAME=Qt6 -DCOMPILER\_ID=GNU \ldots}
  \item Check that all required packages are installed
        (Sections~\ref{app:linux} and~\ref{app:rpi-station}).
  \item If \code{CMAKE\_PREFIX\_PATH} is needed, pass it explicitly:
        \code{cmake -S . -B build -DCMAKE\_PREFIX\_PATH=/path/to/Qt/6.x.x/...}
\end{enumerate}

\subsection{FairWindSK does not autostart on Raspberry Pi}

\begin{enumerate}
  \item Verify the autostart file was installed:
\begin{lstlisting}
ls /etc/xdg/autostart/fairwindsk-startup.desktop
\end{lstlisting}
  \item Verify the \code{fairwindsk-startup} script is executable
        and in PATH:
\begin{lstlisting}
which fairwindsk-startup
\end{lstlisting}
  \item Run the startup helper manually from a terminal to check for errors:
\begin{lstlisting}
fairwindsk-startup
\end{lstlisting}
  \item Ensure the desktop session supports XDG autostart
        (the default Raspberry~Pi~OS desktop does).
\end{enumerate}
